\chapter{Intelligent Features of the System}

\vspace{-0.5cm}
\begin{center}
\begin{minipage}{0.85\textwidth}
\small
\textbf{Contents}\\[0.3cm]
\hspace*{1cm} 5.1 \quad Overview \dotfill \pageref{sec:ch5overview}\\
\hspace*{1cm} 5.2 \quad Ingestion Service --- Document Processing Pipeline \dotfill \pageref{sec:ingestion}\\
\hspace*{1cm} 5.3 \quad Embedding Service --- Semantic Vectorisation \dotfill \pageref{sec:embedding}\\
\hspace*{1cm} 5.4 \quad Search Service --- Two-Stage Semantic Retrieval \dotfill \pageref{sec:search}\\
\hspace*{1cm} 5.5 \quad Summarisation Service \dotfill \pageref{sec:summarisation}\\
\hspace*{1cm} 5.6 \quad RAG Q\&A Service \dotfill \pageref{sec:ragservice}\\
\hspace*{1cm} 5.7 \quad Recommendation Service --- Hybrid NCF + CBF \dotfill \pageref{sec:recommender}\\
\hspace*{1cm} 5.8 \quad Ethical Considerations \dotfill \pageref{sec:ethics}\\
\end{minipage}
\end{center}
\vspace{0.5cm}
\newpage

% ===========================
\section{Overview}
\label{sec:ch5overview}

This chapter provides a comprehensive technical description of the six AI
subsystems embedded in the Learnova platform: the Ingestion Service, the
Embedding Service, the Search Service, the Summarisation Service, the
Retrieval-Augmented Generation (RAG) Q\&A Service, and the Recommendation
Service. Together, these services transform a static collection of book files
into a dynamic, queryable knowledge base that users can explore through natural
language, receive personalised recommendations, and query for specific
information --- without any manual indexing or metadata entry.

% ===========================
\section{Ingestion Service --- Document Processing Pipeline}
\label{sec:ingestion}

The Ingestion Service is the entry point for all book content in the platform.
It accepts PDF and ePub files uploaded by library administrators and
orchestrates a three-stage pipeline: text extraction, cleaning and sentence
tokenisation, and sliding-window chunking. The output is a stream of text
chunks published to Kafka, which the Embedding Service then consumes
asynchronously to produce vector representations.

\subsection{Text Extraction}

Format detection is based on file extension. PDF files are processed using
\textbf{PyMuPDF} (\texttt{fitz}), which extracts plain text page by page.
ePub files are processed using \textbf{ebooklib}, which iterates over HTML
chapter items, strips all HTML tags, and yields plain text sections. In both
cases, pages or sections with fewer than 50 characters are silently skipped,
filtering out blank pages, cover images, and decorative separators that would
otherwise produce noise in the vector index.

\subsection{Cleaning and Chunking}

Extracted text is cleaned and sentence-tokenised before chunking. The chunker
uses the \texttt{all-mpnet-base-v2} tokenizer to measure token counts
precisely, ensuring that chunks match the embedding model's input space. A
sliding-window algorithm fills each chunk greedily up to 512 tokens, then
carries the last 50 tokens into the next chunk as overlap --- preserving
cross-boundary semantic continuity. Sentences that individually exceed 512
tokens are split into sub-windows before chunking begins.

Each chunk is assigned a deterministic identifier of the form
\texttt{\{isbn\}\_p\{page\_num\}\_c\{chunk\_index\}}, enabling exact
traceability from any retrieved vector back to its source page in the original
book.

\begin{table}[H]
\centering
\caption{Ingestion Service Chunking Parameters}
\label{tab:chunking}
\begin{tabularx}{\textwidth}{|l|l|l|X|}
\hline
\textbf{Parameter} & \textbf{Value} & \textbf{Config Key} & \textbf{Effect} \\
\hline
Chunk Size        & 512 tokens & CHUNK\_SIZE        & Maximum tokens per chunk \\
\hline
Chunk Overlap     & 50 tokens  & CHUNK\_OVERLAP     & Tokens carried into next chunk \\
\hline
Min Page Length   & 50 chars   & (hardcoded)        & Filters blank/decorative pages \\
\hline
Upsert Batch Size & 100 vectors & UPSERT\_BATCH\_SIZE & Vectors flushed per Pinecone call \\
\hline
\end{tabularx}
\end{table}

\subsection{Kafka Publishing and Status Tracking}

Each chunk is serialised as a JSON message and published to the Kafka topic
\texttt{book.chunks} with the book's ISBN as the partition key, ensuring all
chunks from the same book are processed in order by the same Embedding Service
consumer instance. Ingestion runs as a FastAPI background task: the API accepts
the upload, stores the file to disk, and immediately returns HTTP 202 Accepted
with a job ID. The caller can poll \texttt{/ingest/\{isbn\}/status} to track
progress through the \texttt{queued} $\rightarrow$ \texttt{processing}
$\rightarrow$ \texttt{done} lifecycle, with status persisted in Redis.

% ===========================
\section{Embedding Service --- Semantic Vectorisation}
\label{sec:embedding}

The Embedding Service is the foundational layer of the platform's AI pipeline.
It converts raw text chunks, produced by the Ingestion Service, into dense
768-dimensional numerical vectors that encode semantic meaning. These vectors
are stored in Pinecone and serve as the retrieval backbone for every downstream
AI capability: semantic search, recommendations, summarisation, and RAG Q\&A.

\subsection{Design Rationale}

Traditional library search systems rely on keyword matching over metadata
fields such as title, author, and ISBN. This approach fails to surface
semantically relevant content buried within the book text itself. Semantic
search addresses this limitation by representing every passage as a point in
a shared vector space: two passages with similar meaning are geometrically
close, even if they share no words in common.

The Embedding Service is designed as a headless Kafka consumer with no HTTP
interface of its own. This architectural choice decouples embedding from the
upload API, allowing large books to be indexed asynchronously without blocking
the user-facing request. Multiple consumer instances can be deployed
independently to scale ingestion throughput.

\subsection{Model: all-mpnet-base-v2}

The \texttt{sentence-transformers/all-mpnet-base-v2} model was selected as the
embedding backbone after evaluation on the Semantic Textual Similarity (STS)
benchmark. It is a fine-tuned variant of Microsoft's MPNet architecture,
trained on over one billion sentence pairs, producing 768-dimensional vectors
with state-of-the-art performance on passage-level retrieval tasks.

A critical mathematical property is L2 normalisation: all output vectors are
projected onto the unit hypersphere. Under this normalisation, cosine
similarity equals the dot product of two vectors --- a property exploited by
Pinecone's ANN index through maximum inner-product search, enabling
sub-millisecond retrieval across millions of vectors.

\begin{table}[H]
\centering
\caption{Embedding Model Configuration and Design Decisions}
\label{tab:embedding}
\begin{tabularx}{\textwidth}{|l|l|X|}
\hline
\textbf{Property} & \textbf{Value} & \textbf{Justification} \\
\hline
Model            & all-mpnet-base-v2       & Top STS benchmark performance on passage retrieval \\
\hline
Output Dimension & 768-d float32           & Balance between representational power and storage cost \\
\hline
Normalisation    & L2 (unit sphere)        & Cosine similarity = dot product, exploited by Pinecone ANN \\
\hline
Batch Size       & 64 texts / forward pass & Optimal GPU utilisation without out-of-memory risk \\
\hline
Hardware         & CUDA GPU (CPU fallback) & Automatic device selection at runtime \\
\hline
Lazy Loading     & On first encode() call  & Avoids GPU initialisation overhead at service startup \\
\hline
\end{tabularx}
\end{table}

\subsection{BookEncoder Implementation}

The \texttt{BookEncoder} class wraps the SentenceTransformer library and
exposes two methods: \texttt{encode\_single()} for real-time query encoding
and \texttt{encode\_batch()} for bulk ingestion. The
\texttt{torch.no\_grad()} context manager disables gradient computation during
inference, reducing memory usage and improving throughput by approximately
30\%.

\begin{lstlisting}[language=Python, caption=BookEncoder encode\_batch method, 
  basicstyle=\ttfamily\footnotesize, breaklines=true]
def encode_batch(self, texts: list[str], batch_size: int = 64) -> np.ndarray:
    model = self._ensure_model()       # lazy-load on first call
    with torch.no_grad():              # disable gradients for inference
        vecs = model.encode(
            texts, batch_size=batch_size,
            convert_to_numpy=True,
            normalize_embeddings=True, # apply L2 normalisation
        )
    return vecs.astype(np.float32)     # 768-d float32 matrix
\end{lstlisting}

\subsection{Kafka-to-Pinecone Pipeline and Fault Tolerance}

The \texttt{EmbeddingBatchJob} class buffers up to 100 vector records in
memory and flushes them to Pinecone as a single batched upsert, reducing API
call overhead. Each vector is stored in a Pinecone namespace equal to the
book's ISBN, enabling efficient per-book filtering without maintaining
separate indexes.

Upsert operations are retried up to three times with exponential back-off
(delays of 1s, 2s, and 4s). If all three attempts fail, the batch is
forwarded to a Kafka dead-letter topic (\texttt{book.chunks.dlq}) for manual
inspection and replay.

\begin{table}[H]
\centering
\caption{Embedding Batch Job Processing Stages}
\label{tab:embedbatch}
\begin{tabularx}{\textwidth}{|l|l|l|X|}
\hline
\textbf{Stage} & \textbf{Input} & \textbf{Output} & \textbf{Technology} \\
\hline
Consume & Kafka: book.chunks      & ChunkData object         & aiokafka KafkaConsumer \\
\hline
Encode  & Chunk text (string)     & 768-d float32 vector     & SentenceTransformer \\
\hline
Buffer  & Vector record + metadata & In-memory list (max 100) & Python list \\
\hline
Upsert  & List of records         & Pinecone index entry     & Pinecone client upsert() \\
\hline
DLQ     & Failed batch            & Kafka: book.chunks.dlq   & aiokafka KafkaProducer \\
\hline
\end{tabularx}
\end{table}

% ===========================
\section{Search Service --- Two-Stage Semantic Retrieval}
\label{sec:search}

The Search Service provides natural-language book search across the entire
library catalogue. Unlike traditional keyword search, the service encodes the
query into a 768-dimensional vector and retrieves semantically similar passages
--- surfacing relevant content even when the query shares no exact words with
the text. To balance recall and precision, the service implements a two-stage
retrieval pipeline combining a fast bi-encoder for candidate recall with a
slower but more accurate cross-encoder for re-ranking.

\subsection{Two-Stage Pipeline}

\textbf{Stage 1 --- Approximate Nearest Neighbour (ANN) Retrieval:} The user's
query is encoded by the same \texttt{all-mpnet-base-v2} bi-encoder used during
ingestion, producing an L2-normalised 768-dimensional vector. This vector is
submitted to Pinecone's global namespace, retrieving the top 50 candidates
(\texttt{ANN\_TOP\_K = 50}) by cosine similarity. An optional genre filter can
be applied as a Pinecone metadata filter to restrict results to specific
categories.

\textbf{Stage 2 --- Cross-Encoder Re-ranking:} The 50 ANN candidates are
passed to a CrossEncoder model
(\texttt{cross-encoder/ms-marco-MiniLM-L-6-v2}), which scores each
query--passage pair jointly. Unlike the bi-encoder, which encodes query and
passage independently, the cross-encoder attends to both simultaneously ---
producing more accurate relevance scores at the cost of higher latency. The
top 10 results (\texttt{FINAL\_TOP\_K = 10}) by cross-encoder score are
returned to the caller.

\begin{table}[H]
\centering
\caption{Search Service Two-Stage Pipeline}
\label{tab:searchpipeline}
\begin{tabularx}{\textwidth}{|l|l|X|}
\hline
\textbf{Stage} & \textbf{Model} & \textbf{Output} \\
\hline
ANN Retrieval  & all-mpnet-base-v2 (bi-encoder) & Top 50 candidates from Pinecone global namespace \\
\hline
Re-ranking     & cross-encoder/ms-marco-MiniLM-L-6-v2 & Top 10 results sorted by cross-encoder score \\
\hline
Similar Books  & CBF (Pinecone synopsis namespace) & Books semantically similar to a given ISBN \\
\hline
\end{tabularx}
\end{table}

\subsection{Caching and Similar-Book Search}

Search results are cached in Redis for 60 seconds using a SHA-256 hash of the
query, top-k, and genre filter as the cache key. This avoids redundant
Pinecone queries for repeated searches within the same minute. The service also
exposes a \texttt{/search/similar/\{isbn\}} endpoint that delegates to the
Content-Based Filtering layer of the Recommendation Service, returning books
whose synopsis vectors are closest to the requested ISBN in Pinecone's
synopses namespace.

% ===========================
\section{Summarisation Service}
\label{sec:summarisation}

The Summarisation Service generates concise textual summaries of full book
content and structured mind-map representations of key themes. Because a full
book may contain thousands of text chunks --- far exceeding any language model
context window --- the service employs a map-reduce strategy to handle
arbitrarily long documents.

\subsection{Model: FLAN-T5-large}

The service uses Google's \textbf{FLAN-T5-large}, a sequence-to-sequence model
based on the T5 architecture. FLAN-T5 was instruction-fine-tuned on a large
collection of tasks expressed as natural language instructions, enabling it to
follow structured prompts reliably without task-specific fine-tuning --- making
it well-suited for a library system where content varies widely in genre and
domain.

Beam search (\texttt{num\_beams=4}) explores multiple candidate sequences
simultaneously, consistently producing higher-quality summaries than greedy
decoding. A \texttt{no\_repeat\_ngram\_size} of 3 suppresses repetitive phrase
patterns, a common failure mode in extractive-style models.

\begin{table}[H]
\centering
\caption{FLAN-T5 Prompt Configurations by Summary Type}
\label{tab:summarytypes}
\begin{tabularx}{\textwidth}{|l|X|l|l|}
\hline
\textbf{Type} & \textbf{Prompt Role} & \textbf{Max Tokens} & \textbf{Use Case} \\
\hline
SHORT    & Book blurb writer --- 2--3 engaging sentences for a general reader & 150 & Quick overview \\
\hline
DETAILED & Literary analyst --- thematic summary with key events and characters & 300 & Deep-dive analysis \\
\hline
\end{tabularx}
\end{table}

\subsection{Map-Reduce Pipeline}

The \texttt{MapReducePipeline} processes books in two phases to overcome the
context window constraint:

\textbf{MAP Phase:} Chunks are fetched from Pinecone in page-number order and
divided into batches of 5 consecutive chunks. Each batch is independently
summarised by FLAN-T5 with an 80-token output cap, producing a list of
intermediate summaries. Batches are processed sequentially to respect GPU
memory constraints.

\textbf{REDUCE Phase:} All intermediate summaries are concatenated into a
single string and passed through FLAN-T5 one final time with the user-selected
type --- SHORT (150 tokens) or DETAILED (300 tokens) --- producing the final
book summary.

The final summary is cached in Redis with a TTL of 24 hours, so repeated
requests for the same book are served instantly without re-running the
pipeline. A GPU queue counter in Redis limits concurrent summarisation jobs to
a maximum of 4. If the queue is full, the API returns HTTP 503 with a
\texttt{Retry-After} header.

\subsection{Mind Map Generation}

In addition to textual summaries, the service can generate a structured JSON
mind-map tree representing the book's key themes. The \texttt{MindMapBuilder}
class prompts FLAN-T5 to produce a JSON object with a root node (the book
title) and child nodes (main themes), each containing a list of supporting
subpoints. All output is validated against a Pydantic schema
(\texttt{MindMapOutput}). If validation fails, a rule-based fallback generator
guarantees a valid, well-structured response.

\begin{table}[H]
\centering
\caption{Mind Map Generation Pipeline with Fallback Strategy}
\label{tab:mindmap}
\begin{tabularx}{\textwidth}{|l|X|X|l|}\hline
\textbf{Stage} & \textbf{Input} & \textbf{Output} & \textbf{Fallback?} \\
\hline
Prompt              & Book summary (max 3000 chars) & Structured prompt  & No \\
\hline
Generation          & Prompt + FLAN-T5              & Raw text output    & No \\
\hline
JSON Extraction     & Raw text                      & Parsed dict object & Yes --- rule-based \\
\hline
Pydantic Validation & dict object                   & MindMapOutput      & Yes --- rule-based \\
\hline
Response            & MindMapOutput                 & JSON tree          & Fallback always valid \\
\hline
\end{tabularx}
\end{table}

% ===========================
\section{RAG Q\&A Service}
\label{sec:ragservice}

The RAG Service enables users to pose natural-language questions about any book
in the library and receive accurate, page-cited answers. Rather than relying on
a language model's parametric knowledge, the service first retrieves the most
relevant passages from the book and provides them as explicit context to the
LLM, grounding its response in the actual text.

\subsection{Architecture}

The service is built on the \textbf{LangChain} framework with
\textbf{Groq's Llama-3.3-70b-Versatile} as the generation model and Pinecone
as the vector retrieval backend. The pipeline consists of three sequential
components: a retriever, a chain, and a streaming layer.

\begin{table}[H]
\centering
\caption{RAG Service Component Overview}
\label{tab:ragcomponents}
\begin{tabularx}{\textwidth}{|l|l|X|}
\hline
\textbf{Component} & \textbf{Technology} & \textbf{Responsibility} \\
\hline
Retriever & PineconeBookRetriever & Encodes query, fetches top-K passages by ISBN namespace \\
\hline
Chain     & LangChain LCEL + PromptTemplate & Formats passages as context and submits to LLM \\
\hline
LLM       & Llama-3.3-70b (temp = 0.1) & Generates grounded answer citing page numbers via Groq \\
\hline
Streaming & SSE + asyncio Queue & Delivers tokens incrementally as LLM produces them \\
\hline
Cache     & LRU cache (256 entries) & Avoids rebuilding chain on repeated requests \\
\hline
\end{tabularx}
\end{table}
\subsection{Retriever --- Similarity Threshold Filtering}

The \texttt{PineconeBookRetriever} encodes the user's question using the same
\texttt{all-mpnet-base-v2} model used during ingestion --- query and document
vectors must occupy the same semantic space for similarity search to be
meaningful. The resulting 768-dimensional query vector is submitted to Pinecone
for ANN search, retrieving the top 5 candidate passages
(\texttt{RAG\_TOP\_K = 5}) within the namespace corresponding to the requested
book's ISBN.

A similarity threshold of \textbf{0.30} is applied to the retrieved results:
any passage scoring below this threshold is discarded before context assembly.
This mechanism serves as a \textbf{hallucination guard} --- without it, a
model receiving weakly-related passages may generate plausible-sounding answers
not supported by the book. If no passages exceed the threshold, the chain
returns a standardised insufficient-information message.

\subsection{Chain --- Grounded Prompting and Token Budget}

The RAG pipeline is implemented as a LangChain Expression Language (LCEL)
chain following the ``stuff'' strategy, concatenating all retrieved passages
into a single context block. A custom system prompt enforces three strict
behavioural constraints:

\begin{itemize}
  \item Answer using \textbf{only} the provided excerpts --- no external or
        parametric knowledge is permitted.
  \item Cite the page number(s) from which each piece of information was drawn.
  \item Respond with a standardised message if the provided excerpts are
        insufficient to answer the question.
\end{itemize}

\begin{lstlisting}[language=Python, caption=RAG System Prompt Definition,
  basicstyle=\ttfamily\footnotesize, breaklines=true]
SYSTEM_PROMPT = (
  "You are an intelligent reading assistant for the Smart AI Library. "
  "Answer using ONLY the excerpts provided below. "
  "If the excerpts are insufficient, respond: "
  "'I don't have enough information from this book to answer that.' "
  "Always cite the page number(s) you drew information from."
)
\end{lstlisting}

A token budget guard truncates the assembled context to a maximum of 3,800
tokens using \texttt{tiktoken} (\texttt{cl100k\_base} encoding), preventing
context-overflow API failures.

\subsection{Streaming --- Server-Sent Events (SSE)}

To improve perceived responsiveness, the service delivers
Llama-3.3-70b-Versatile's response token-by-token via Server-Sent Events
(SSE). An \texttt{asyncio Queue} acts as an in-process message bus between the
LangChain streaming callback and the SSE generator coroutine. A 15-second
inter-token timeout terminates the stream if the model stalls.

\begin{table}[H]
\centering
\caption{SSE Protocol Events in the RAG Streaming Endpoint}
\label{tab:sseevents}
\begin{tabularx}{\textwidth}{|l|X|}
\hline
\textbf{SSE Event} & \textbf{Meaning} \\
\hline
\texttt{data: <token>}   & A generated token --- client appends it to the displayed answer \\
\hline
\texttt{data: [DONE]}    & Generation complete --- client should close the event stream \\
\hline
\texttt{data: [TIMEOUT]} & No token received within 15 seconds --- stream terminated \\
\hline
\texttt{data: [ERROR]}   & Exception during chain execution --- stream aborted \\
\hline
\end{tabularx}
\end{table}

% ===========================
\section{Recommendation Service --- Hybrid NCF + CBF}
\label{sec:recommender}

The Recommendation Service delivers personalised book recommendations by
combining two complementary AI techniques: Neural Collaborative Filtering
(NCF), which learns from user behaviour, and Content-Based Filtering (CBF),
which reasons over book content semantics. A hybrid fusion layer blends the
two signals with weights that adapt dynamically based on the user's interaction
history.

\subsection{Neural Collaborative Filtering (NCF)}

The NCF model is a PyTorch implementation of the Generalised Matrix
Factorisation + MLP architecture (He et al., 2017). It maintains two separate
embedding tables for users and items --- one for the GMF branch and one for
the MLP branch --- allowing each pathway to learn complementary interaction
patterns.

The \textbf{GMF branch} computes an element-wise product of user and item
embeddings, capturing linear co-occurrence signals. The \textbf{MLP branch}
concatenates the embeddings and passes them through three fully-connected
layers [128, 64, 32] with ReLU activations and 0.2 dropout, capturing
non-linear interaction patterns. The outputs of both branches are concatenated
and passed through a final sigmoid layer to produce a relevance score in
$[0,1]$.

\begin{table}[H]
\centering
\caption{NCF Model Hyperparameters}
\label{tab:ncf}
\begin{tabularx}{\textwidth}{|l|X|}
\hline
\textbf{Hyperparameter} & \textbf{Value} \\
\hline
Embedding Dimension (emb\_dim) & 64 \\
\hline
MLP Hidden Layers & [128, 64, 32] \\
\hline
Dropout Rate & 0.2 \\
\hline
Output Activation & Sigmoid $\rightarrow$ score in $[0, 1]$ \\
\hline
Weight Initialisation & Normal (std=0.01) for embeddings; Xavier uniform for linear layers \\
\hline
\end{tabularx}
\end{table}

The model is updated incrementally via an Online Updater that consumes user
interaction events from the Kafka topic \texttt{user.events}. Events classified
as positive signals (\texttt{read\_start}, \texttt{bookmark},
\texttt{rating\_4}, \texttt{rating\_5}) trigger a single SGD step on the user
embedding layers only (learning rate = $5 \times 10^{-4}$), keeping item
embeddings stable while personalising user representations in near real-time.

\subsection{Content-Based Filtering (CBF)}

The CBF layer uses Pinecone's \texttt{synopses} namespace, which stores one
synopsis vector per book. Each synopsis vector is the L2-normalised average of
all chunk vectors for that book, computed and upserted by the Embedding Service
after ingestion. To find books similar to a seed title, the CBF layer fetches
the seed's synopsis vector and queries Pinecone for the closest neighbours by
cosine similarity.

\subsection{Hybrid Fusion and User Tiers}

The final recommendation score is a weighted linear blend of the NCF and CBF
scores:

\begin{equation}
\text{score} = \alpha \times \text{CF} + (1 - \alpha) \times \text{CBF}
\end{equation}

The weight $\alpha$ is determined by the user's interaction tier, computed
from the number of reading events recorded in the database:

\begin{table}[H]
\centering
\caption{Hybrid Fusion Alpha Values by User Tier}
\label{tab:usertier}
\begin{tabularx}{\textwidth}{|l|l|X|}
\hline
\textbf{Tier} & \textbf{Condition} & \textbf{Alpha ($\alpha$)} \\
\hline
New       & Fewer than 5 interactions  & 0.10 --- mostly content-based \\
\hline
Returning & 5 to 49 interactions       & 0.55 --- balanced blend \\
\hline
Power     & 50 or more interactions    & 0.80 --- mostly collaborative \\
\hline
\end{tabularx}
\end{table}

This tiered approach solves the cold-start problem: new users with little
interaction history receive content-driven recommendations based on book
semantics, while power users benefit primarily from collaborative signals
learned from their accumulated behaviour. Recommendations are cached in Redis
for 5 minutes per user and are subject to a rate limit of 60 requests per
minute per user, enforced via a Redis sliding-window counter.

% ===========================
\section{Ethical Considerations}
\label{sec:ethics}

The deployment of AI-powered features in a library context carries
responsibilities beyond technical correctness. The following ethical principles
were explicitly considered during the design and implementation of the
platform's AI subsystems.

\subsection{Hallucination Mitigation}

Large language models can generate factually incorrect information with high
confidence --- a phenomenon termed \textit{hallucination}. In a library
context, a hallucinated answer could seriously mislead researchers or students.
The RAG service addresses this through two complementary mechanisms: the
similarity threshold filter rejects weakly-related passages before they reach
the LLM, and the system prompt explicitly forbids the model from drawing on
knowledge outside the provided excerpts.

\subsection{Transparency and Source Attribution}

All RAG Q\&A answers include explicit page-number citations from the source
documents returned by the retriever. Users can verify any claim by consulting
the original book at the cited page. This design treats AI-generated answers
as a navigational aid rather than an authoritative endpoint, preserving the
user's ability to critically evaluate information.

\subsection{Content Scope and Copyright}

The platform processes only book files uploaded by authenticated
administrators, who are assumed to have verified copyright compliance before
upload. Generated summaries are abstractive paraphrases intended to assist
comprehension, not to substitute for reading the original work or to
circumvent copyright restrictions.

\subsection{Model Bias and Limitations}

The pre-trained models used in this platform
(\texttt{all-mpnet-base-v2}, \texttt{FLAN-T5-large},
\texttt{Llama-3.3-70b-Versatile}, and
\texttt{cross-encoder/ms-marco-MiniLM-L-6-v2}) were not fine-tuned on
library-domain data. As a result, they may perform less accurately on highly
specialised technical content or non-English texts. The platform currently
supports English-language content; extension to Arabic and other languages is
identified as a planned future enhancement.

% ===========================
\section*{Chapter Summary}
\addcontentsline{toc}{section}{Chapter Summary}

This chapter documented the six core AI subsystems of the Learnova platform.
The Ingestion Service processes PDF and ePub files through a three-stage
pipeline --- text extraction, chunking, and Kafka publishing --- producing
traceable 512-token chunks for downstream services. The Embedding Service
consumes these chunks and encodes them into L2-normalised 768-dimensional
vectors using \texttt{all-mpnet-base-v2}, stored in Pinecone by ISBN
namespace. The Search Service combines a bi-encoder ANN stage with a
cross-encoder re-ranking stage to deliver accurate semantic search results.
The Summarisation Service applies a FLAN-T5-large map-reduce pipeline to
produce SHORT and DETAILED summaries and structured mind-map trees. The RAG
Q\&A Service combines Pinecone retrieval with Llama-3.3-70b-Versatile
generation, enforcing grounded output through a similarity threshold and a
strict system prompt, delivering responses via SSE streaming with page-level
citations. Finally, the Recommendation Service blends NCF and CBF signals
through a tiered fusion strategy, adapting dynamically to each user's
interaction history. Together, these services transform Learnova from a static
digital library into an intelligent, interactive knowledge platform.